\documentclass[12pt,a4paper]{ctexart}
% 公共排版文件（使用 XeLaTeX 编译）
\usepackage[a4paper,margin=2.35cm]{geometry}
\usepackage{booktabs,longtable,tabularx,array,multirow,enumitem}
\usepackage{graphicx,xcolor,amsmath,hyperref,lastpage,fancyhdr}
\usepackage{tikz}
\usetikzlibrary{arrows.meta,positioning,shapes.geometric,fit,calc}
\hypersetup{colorlinks=true,linkcolor=blue!60!black,urlcolor=blue!60!black,citecolor=blue!60!black,pdftitle={逸仙活动云课程报告}}
\setlist{itemsep=2pt,topsep=3pt,leftmargin=2em}
\setlength{\parindent}{2em}
\setlength{\headheight}{15pt}
\renewcommand{\arraystretch}{1.35}
\newcolumntype{Y}{>{\raggedright\arraybackslash}X}
\newcommand{\reportcover}[1]{%
  \begin{titlepage}
  \begin{center}
  \vspace*{3cm}{\zihao{1}\bfseries 中山大学\\[0.6cm]《软件工程》期末大作业\par}
  \vspace{2.5cm}{\zihao{2}\bfseries #1\par}
  \end{center}
  \vspace{2cm}
  \begin{center}
  \begin{tabular}{l}
  {\Large 项目名称：逸仙活动云——中山大学校园活动云平台}\\
  {\Large 小组成员：钟梓俊\quad 邱剑波\quad 张宸睿}\\
  {\Large 指导老师：王可泽}
  \end{tabular}
  \end{center}
  \vfill
  \begin{center}
  {\large 2026 年 7 月 16 日\par}
  \end{center}
  \end{titlepage}}

\pagestyle{fancy}\fancyhf{}\lhead{逸仙活动云}\rhead{软件工程期末大作业}\cfoot{第 \thepage\ 页，共 \pageref{LastPage} 页}
\setcounter{tocdepth}{2}

\begin{document}
\reportcover{架构设计文档}
\pagenumbering{Roman}\tableofcontents\clearpage\pagenumbering{arabic}
\docmeta

\section{架构设计概述}
\subsection{设计目标}
系统架构围绕四个核心目标展开：
\begin{enumerate}
\item \textbf{可演示：} 支持在本地开发环境（SQLite + Mock）完整运行核心业务流程，不依赖外部邮件、搜索或 AI 服务即可展示活动浏览、认证、发布审核和个人参与等主要功能。
\item \textbf{可联调：} 前后端通过预定义的 REST/JSON 接口契约并行开发，Mock 服务允许前端在无后端环境时独立验证界面逻辑。
\item \textbf{可部署：} 提供 Docker Compose 容器化部署拓扑，单命令即可在 Linux 服务器上重建完整运行环境。
\item \textbf{可演化：} 分层架构和服务层隔离确保业务规则变更时影响范围可控，领域名称适配层隔离前后端数据结构差异。
\end{enumerate}

\subsection{设计原则}
\begin{enumerate}
\item \textbf{接口优先：} 前后端以接口契约为核心协作依据，字段、枚举、分页和时区等约定在编码前达成一致。
\item \textbf{依赖反转：} 前端依赖抽象接口而非具体实现，后端服务层与基础设施层通过依赖注入解耦。
\item \textbf{最小权限：} 每个 API 端点服务端校验 JWT 身份与 RBAC 角色，前端不单独承担授权职责。
\item \textbf{可观测性：} 关键操作写入审计日志，搜索日志脱敏记录，异常统一格式化为 JSON 响应。
\item \textbf{配置外置：} 密钥、数据库连接、外部服务地址通过环境变量注入，不硬编码进入源码。
\end{enumerate}

\subsection{设计约束}
\begin{enumerate}
\item \textbf{技术栈：} 前端使用 Vue 3 + TypeScript + Vite + Pinia，后端使用 Flask + SQLAlchemy + JWT，数据库使用 PostgreSQL（开发环境可用 SQLite 替代）。
\item \textbf{部署环境：} 课程项目以单机可演示为首要目标，Docker Compose 定义完整服务拓扑；生产环境需额外补充反向代理、TLS 和监控。
\item \textbf{外部依赖：} 邮件发送、外部搜索和 AI 增强为可选能力，设计降级方案确保核心流程不因外部不可用而中断。
\item \textbf{开发周期：} 项目周期约 12 周，前端与后端并行开发，通过 Mock 服务和接口契约解耦联调依赖。
\end{enumerate}

\section{系统总体架构}
\subsection{总体架构图}
\begin{center}
\begin{tikzpicture}[font=\small,
  box/.style={draw,rounded corners,align=center,minimum width=2.6cm,minimum height=.75cm,fill=blue!7},
  svc/.style={draw,rounded corners,align=center,minimum width=2.3cm,minimum height=.65cm,fill=green!8},
  arr/.style={-Latex,thick}
]
\node[box] (browser) at (0,2.8) {浏览器\\访客/用户/发布者/管理员};
\node[box] (fe) at (4,2.8) {Vue 3 前端\\Vite · Router · Pinia};
\node[box] (api) at (8,2.8) {Flask API\\蓝图 · JWT · 服务层};
\node[svc] (db) at (11.8,4.0) {PostgreSQL/SQLite\\SQLAlchemy};
\node[svc] (redis) at (11.8,2.5) {Redis\\缓存/限流/队列};
\node[svc] (worker) at (11.8,1.0) {Celery Worker/Beat\\抓取、AI、定时任务};
\node[svc] (external) at (8,.4) {邮件、搜索、AI、外部数据源};
\draw[arr] (browser)--node[above]{HTTPS}(fe);
\draw[arr] (fe)--node[above]{/api, JSON, Bearer JWT}(api);
\draw[arr] (api)--(db);  \draw[arr] (api)--(redis);
\draw[arr] (worker)--(db);  \draw[arr] (worker)--(redis);
\draw[arr,dashed] (api)--node[right]{enqueue}(redis);
\draw[arr,dashed] (redis)--node[right]{dequeue}(worker);
\draw[arr] (api)--(external);  \draw[arr] (worker)--(external);
\end{tikzpicture}
\vspace{0.2em}
{\small\itshape 图1：系统总体架构——前后端分离，异步任务通过 Redis 队列解耦}
\end{center}
前端开发阶段中，Vite 将 \texttt{/api} 请求代理到后端开发服务器；生产环境则由反向代理统一接管静态资源分发、API 路由和 TLS 终止。API 服务将耗时任务（抓取、AI 处理、定时通知）封装为消息放入 Redis 队列，Celery Worker 异步消费，避免阻塞 Web 请求。

\subsection{架构模式}
系统采用三种核心架构模式：
\begin{enumerate}
\item \textbf{前后端分离模式：} 前端为 Vue 3 单页应用，后端为 Flask REST API，通过 HTTP JSON 通信。两者独立构建、独立部署，通过接口契约保证一致性。前端在开发期可使用 Mock 服务先行开发，不等待后端接口就绪。
\item \textbf{分层架构模式：} 后端按职责划分为接口层、服务层、持久层和基础设施层，每层只与相邻层交互。服务层封装可复用的业务规则，不依赖 HTTP 上下文，便于单元测试。
\item \textbf{异步任务模式：} 对于抓取、AI 增强等耗时操作，API 将任务描述写入 Redis 队列后即刻返回，Celery Worker 异步执行并写回结果。任务状态和失败重试由 Worker 管理，不阻塞核心 API 路径。
\end{enumerate}

\subsection{系统分层}
系统自顶向下分为五个层次：
\begin{enumerate}
\item \textbf{表现层：} Vue 3 前端，包含页面（views）、可复用组件（components）、路由（router）和状态管理（stores）。
\item \textbf{接口层：} Flask 蓝图，负责 HTTP 路由、参数解析、身份校验和响应序列化。
\item \textbf{服务层：} 封装业务规则和质量策略，包括内容解析、质量评分、去重、安全校验、审计和通知。
\item \textbf{持久层：} SQLAlchemy 模型，定义实体关系、约束条件、事务管理和数据序列化。
\item \textbf{基础设施层：} 数据库（PostgreSQL/SQLite）、缓存与队列（Redis）、异步任务（Celery）、以及外部服务（邮件、搜索、AI）的客户端适配。
\end{enumerate}

\subsection{各层职责}
\begin{tabularx}{\textwidth}{p{2.3cm}Y Y}
\toprule 层次 & 职责 & 代表实现 \\\midrule
表现层 & 页面组装、路由分发、状态管理、接口适配 & views、components、stores、api 模块 \\
接口层 & HTTP 路由、参数解析、JWT 校验、角色鉴权 & auth、activities、search 等 Flask 蓝图 \\
服务层 & 业务规则、内容质量、去重、审计、通知、安全校验 & \texttt{services/*} \\
持久层 & 实体定义、关系映射、约束、事务 & User、Poster、Registration、AuditLog 等模型 \\
基础设施层 & 数据库、缓存、队列、外部服务适配 & extensions、config、Docker Compose \\
\bottomrule
\end{tabularx}

\section{功能模块设计}
\subsection{模块划分}
依据需求分析的功能域划分，系统包含以下七个核心功能模块和一个通用模块：
\begin{enumerate}
\item \textbf{活动发现模块：} 活动列表、关键词搜索、分类/状态筛选、活动详情展示。
\item \textbf{身份认证模块：} 图形验证码登录、邮箱注册、JWT 会话管理、角色鉴权。
\item \textbf{个人参与模块：} 报名、取消报名、收藏、个人日历、ICS 导出。
\item \textbf{活动发布模块：} 草稿创建与编辑、提交审核、审核结果查看。
\item \textbf{审核管理模块：} 活动审核、批量操作、发布者申请审批、审核意见记录。
\item \textbf{知识与通知模块：} 知识节点关联、订阅管理、通知生成与已读标记。
\item \textbf{数据治理模块：} 数据源维护、抓取日志、内容质量分、疑似重复标识、审计日志。
\item \textbf{异常与通用模块：} 跨页面加载态、空态、错误态、401/403 统一处理。
\end{enumerate}

\subsection{模块职责}
各模块在后端由独立的 Flask 蓝图承载，前端由对应的视图页面和 API 封装模块支撑：
\begin{tabularx}{\textwidth}{p{2.6cm}p{3.4cm}Y}
\toprule 模块 & 后端蓝图/服务 & 前端视图/组件 \\\midrule
活动发现 & posters 蓝图、search 蓝图 & HomeView、SearchView、ActivityDetail \\
身份认证 & auth 蓝图 & LoginView、RegisterView、CaptchaField \\
个人参与 & posters（报名/收藏端点）、calendar 蓝图 & ProfileView、CalendarView \\
活动发布 & posters 蓝图、poster\_service & ActivityEditor、ActivityManagement \\
审核管理 & review 蓝图、audit\_service & ReviewView、AdminDashboard \\
知识与通知 & knowledge 蓝图、notification\_service & KnowledgeManagement \\
数据治理 & data\_sources 蓝图、quality\_service、dedup\_service & DataSourcesView \\
通用 & 全局异常处理、请求日志 & PageState、AppShell \\
\bottomrule
\end{tabularx}

\subsection{模块依赖关系}
模块间的依赖遵循以下规则：
\begin{enumerate}
\item 身份认证模块被所有需登录的模块依赖，提供 JWT 校验和角色查询能力。
\item 活动发现模块无上游依赖，访客和未登录用户均可访问。
\item 活动发布模块依赖于身份认证模块（发布者角色校验），输出待审核活动供审核管理模块消费。
\item 审核管理模块在审核通过后触发知识与通知模块，生成关联关系和站内通知。
\item 数据治理模块作为管理后台能力，独立于用户侧流程，仅管理员可访问。
\item 通用模块被所有页面依赖，提供一致的异常处理、加载态和空态展示。
\end{enumerate}

\section{数据与接口设计}
\subsection{数据存储方案}
系统采用混合存储策略：
\begin{enumerate}
\item \textbf{关系型数据库（PostgreSQL / SQLite）：} 存储所有业务实体，包括用户、活动、报名、收藏、日历事件、知识节点、数据源配置和审计日志。开发环境使用 SQLite 实现零配置启动，部署环境切换为 PostgreSQL。两者通过 SQLAlchemy ORM 抽象层无缝切换，业务代码无需感知底层数据库差异。
\item \textbf{Redis：} 承载三方面职责——缓存（活动列表热点数据）、限流（登录和 API 频率控制）、消息队列（Celery 任务调度）。Redis 为可选依赖，不可用时核心业务仍可运行，仅限流和异步任务降级。
\item \textbf{文件存储：} 用户头像和活动附件存储在本地文件系统，开发环境使用项目目录下的上传目录，生产环境建议切换为对象存储。
\end{enumerate}

\subsection{核心数据关系}
系统核心实体关系如下：
\begin{enumerate}
\item \textbf{User—Activity（1:N）：} 一个用户可创建多个活动，活动必须有创建者。
\item \textbf{User—Registration—Activity（M:N）：} 用户与活动之间的报名关系，\texttt{(user\_id, activity\_id)} 唯一约束保证幂等。
\item \textbf{User—Favorite—Activity（M:N）：} 收藏关系，同样由唯一约束保证幂等。
\item \textbf{User—CalendarEvent—Activity（M:N）：} 个人日历事件，支持 ICS 导出。
\item \textbf{Activity—KnowledgeNode（M:N）：} 通过中间表 \texttt{PosterNode} 关联，\texttt{(poster\_id, node\_id, relation\_type)} 唯一。
\item \textbf{User—AuditLog（1:N）：} 关键管理操作必须记录审计日志，包含操作者、动作、目标和时间。
\end{enumerate}

\subsection{模块接口}
模块间以 RESTful API 为通信契约：
\begin{enumerate}
\item \textbf{接口风格：} 所有接口遵循 REST 风格，使用 HTTP 方法语义（GET 查询、POST 创建、PUT 更新、DELETE 删除），统一返回 JSON 格式。
\item \textbf{认证方式：} 接口通过 JWT Bearer Token 认证，Token 在登录时签发，前端通过 Axios 拦截器自动附加到每个请求的 \texttt{Authorization} 头。
\item \textbf{统一响应结构：} 成功响应直接返回数据体；错误响应统一为 \texttt{\{"error": "描述信息"\}} 格式，HTTP 状态码与语义一致（400 参数错误、401 未认证、403 越权、404 资源不存在、500 服务端错误）。
\item \textbf{分页约定：} 列表接口返回 \texttt{\{items: [...], page: number, per\_page: number, total: number\}}，默认每页 20 条。
\item \textbf{时间格式：} 所有时间字段使用 ISO 8601 格式（\texttt{YYYY-MM-DDTHH:mm:ss}），前后端统一解析。
\end{enumerate}

\subsection{外部接口}
系统与三类外部服务交互，均为可选依赖：
\begin{enumerate}
\item \textbf{邮件服务：} 用于发送验证码和通知邮件。通过环境变量配置 SMTP 参数；服务不可达时注册流程降级为提示用户稍后重试，不影响已登录用户的操作。
\item \textbf{搜索服务：} 提供全文检索能力。内部搜索基于数据库 LIKE 查询为默认实现，外部搜索引擎（如 Elasticsearch）为可选增强；外部搜索超时 10s，超时后回退到内部搜索。
\item \textbf{AI 服务：} 用于活动摘要生成和内容推荐。通过 REST API 调用 LLM 服务，超时 30s，超时或不可用时返回基础摘要或默认推荐列表，不影响活动浏览和审核核心路径。
\end{enumerate}

\section{关键设计决策}
\subsection{技术选型}
\begin{longtable}{p{2.2cm}p{3.0cm}p{3.0cm}p{3.0cm}}
\toprule 决策 & 选型方案 & 理由 & 替代方案 \\\midrule\endhead
前端框架 & Vue 3 + TypeScript & 类型安全、组合式 API 逻辑复用、社区活跃 & React、Angular \\
状态管理 & Pinia & 官方推荐、TypeScript 友好、API 简洁 & Vuex、Zustand \\
后端框架 & Flask + 蓝图 & 轻量灵活、生态成熟、学习成本低 & Django、FastAPI \\
ORM & SQLAlchemy & 支持多数据库（SQLite/PostgreSQL）、迁移工具链完善 & Peewee、Django ORM \\
认证 & JWT + RBAC & 无状态、跨域友好、角色边界明确 & Session/Cookie、OAuth \\
容器化 & Docker + Compose & 服务依赖可复现、单命令部署 & Kubernetes、手动部署 \\
\bottomrule
\end{longtable}

\subsection{关键方案说明}
\begin{enumerate}
\item \textbf{RBAC 权限模型：} 系统定义三种角色——\texttt{viewer}（默认）、\texttt{publisher}（发布者）、\texttt{admin}（管理员）。后端在每个受限接口通过装饰器校验 JWT 中的角色声明，前端路由守卫仅做导航引导，不替代服务端授权。
\item \textbf{活动审核状态机：} 活动状态转换由服务层集中控制，包含四种状态——\texttt{draft}（草稿）$\rightarrow$ \texttt{pending\_review}（待审）$\rightarrow$ \texttt{published}（已发布）$\mid$ \texttt{rejected}（已驳回），驳回后可重返 \texttt{draft} 修改。非法状态转换返回 400 错误。
\item \textbf{领域对象适配：} 前端以 \texttt{Activity} 为领域概念，后端持久化模型中为 \texttt{Poster}。前端 API 层负责字段重命名和格式转换，页面组件不直接引用后端模型字段，隔离数据库重构风险。
\end{enumerate}

\subsection{方案权衡}
\noindent\begin{tabularx}{\textwidth}{p{1.3cm}p{2.8cm}p{3.0cm}p{3.2cm}}
\toprule 序号 & 权衡点 & 所选方案的收益 & 付出的代价 \\\midrule
1 & 前后端分离 vs 一体式开发 & 解耦迭代、Mock 先行、独立部署 & 接口协调成本增加，需契约矩阵和联调清单管理 \\
2 & JWT 无状态 vs Session 有状态 & 无需服务端存储会话、跨域友好 & Token 吊销需黑名单机制，生产需配合短生命周期和刷新 Token \\
3 & 关系型 vs NoSQL & 事务一致性、复杂查询支持、唯一约束兜底幂等 & Schema 变更需迁移工具，不适合超高并发写入 \\
4 & Docker Compose vs 裸机部署 & 环境一致、依赖可复现、学习成本低 & 非高可用方案，生产需演化到编排平台（K8s）和托管数据库 \\
\bottomrule
\end{tabularx}

\section{架构质量属性}
\subsection{可维护性}
\begin{enumerate}
\item \textbf{分层隔离：} 前端按 views / components / api / stores / composables 分层，后端按接口 / 服务 / 持久层分层，每层职责明确，修改范围可控。
\item \textbf{领域名称适配：} 前端 API 层统一将后端的 \texttt{Poster} 适配为前端的 \texttt{Activity}，页面代码不直接依赖后端持久化字段。后端表结构重构时，仅需修改 API 适配层，前端组件不受影响。
\item \textbf{测试覆盖：} 服务层函数为纯业务逻辑，不依赖 HTTP 上下文，可独立编写单元测试。后端服务层和前端业务工具函数测试覆盖率目标不低于 70\%，核心模块（审核、权限、质量）不低于 85\%。
\item \textbf{配置管理：} 所有环境差异项通过 \texttt{.env} 文件管理，仓库只保留 \texttt{.env.example} 模板。密钥、数据库连接串和外部服务地址均不进入版本控制。
\end{enumerate}

\subsection{可扩展性}
\begin{enumerate}
\item \textbf{模块化扩展：} 后端使用 Flask 蓝图组织路由，新增功能域只需注册新蓝图并声明路由前缀，不影响已有模块。前端对应新增视图页面和 API 封装模块即可。
\item \textbf{外部能力可插拔：} 邮件、搜索、AI 等服务均通过环境变量开关控制启用状态，服务不可达时系统以降级模式运行。新增外部服务只需实现适配接口并注册到服务层。
\item \textbf{数据源扩展：} 数据源抓取采用插件化设计，新增抓取目标只需配置 URL、允许域和抓取模式，无需修改抓取引擎核心逻辑。
\end{enumerate}

\subsection{性能与可靠性}
\begin{enumerate}
\item \textbf{异步处理：} Celery Worker 承担抓取、AI 处理、定时通知等耗时任务，API 服务通过 Redis 队列异步派发，避免 Web 请求阻塞。Worker 崩溃不影响 API 可用性，任务在 Redis 中持久化等待重新调度。
\item \textbf{超时与降级：} 外部调用设置超时阈值——搜索 10s、AI 服务 30s。超时后返回降级结果（基础搜索或默认推荐），核心活动浏览、认证和审核流程不因外部超时而阻塞。
\item \textbf{缓存策略：} 热点活动列表和可枚举字典数据通过 Redis 缓存，减少数据库查询频率。缓存过期时间为 5--30 分钟，视数据更新频率调整。
\item \textbf{安全防护：} 认证端点实施验证码和限流防止暴力破解；密码以哈希形式存储；外部 URL 抓取校验协议、域名和回环地址以降低 SSRF 风险；所有审核和管理操作记录审计日志以供追溯。
\item \textbf{数据可靠性：} 报名、收藏等关系操作依赖数据库唯一约束提供最终幂等保证。写操作遵循"校验—领域处理—事务提交"的顺序，单次请求内的状态更新、审计记录和通知触发在同一事务边界内完成。
\end{enumerate}

\end{document}
